% chapters/requirements/rf_subjects.tex
% RF-4: Asignaturas y grupos
% ============================================================================

\item \textbf{Asignaturas y grupos.} \label{rf:subjects}
Las asignaturas constituyen la unidad organizativa principal del dominio académico. Cada
asignatura pertenece a un curso académico (que a su vez pertenece a una organización) y
contiene grupos de estudiantes, profesores asignados y exactamente un coordinador. Los
grupos son exclusivos de cada asignatura: un identificador de grupo en una asignatura es
independiente del mismo identificador en otra. Un usuario puede pertenecer a varias
asignaturas del mismo curso y su rol se define en cada una mediante
\texttt{SubjectMembership}.

\begin{functional}

    % RF-4.1 ----------------------------------------------------------------
    \item \textbf{Creación individual de asignatura.} \label{rf:subjects:create}
    El sistema debe permitir a un gestor crear una asignatura asociada al curso activo de
    su organización.

    \textbf{Entrada:} petición POST al recurso de asignaturas con los datos: nombre,
    código, cuatrimestre (etiqueta libre: <<1.er cuatrimestre>>, <<anual>>, etc.) y el
    identificador del usuario que será coordinador.

    \textbf{Proceso:} el sistema valida que el código de asignatura sea único dentro del
    curso activo y que el usuario designado como coordinador exista, esté activo y
    pertenezca a la misma organización. Se crea la asignatura y se asigna al coordinador
    mediante una \texttt{SubjectMembership} con \texttt{role=COORDINATOR} y el conjunto
    completo de permisos contextuales por defecto.

    \textbf{Salida:} respuesta con los datos de la asignatura creada y código
    \acs{http}~201.

    % RF-4.2 ----------------------------------------------------------------
    \item \textbf{Importación masiva de asignaturas.} \label{rf:subjects:import}
    El sistema debe permitir a un gestor crear múltiples asignaturas simultáneamente,
    incluyendo la asignación de usuarios y grupos.

    \textbf{Entrada:} petición POST con un fichero en formato \acs{csv} o \acs{json} que
    contenga, para cada asignatura: nombre, código, cuatrimestre, coordinador, lista de
    profesores, lista de grupos y asignación de estudiantes a grupos.

    \textbf{Proceso:} el sistema parsea el fichero y procesa cada asignatura de forma
    transaccional. Para cada una, se validan los datos, se crea la asignatura, se asignan
    los grupos, se asigna al coordinador y se vinculan profesores y estudiantes con sus
    roles y grupos respectivos mediante \texttt{SubjectMembership}
    (\texttt{is\_active=True}). Si un usuario referenciado no existe en la organización, se
    reporta como error en esa fila sin abortar el resto de la importación.

    \textbf{Salida:} respuesta con un informe del resultado: número de asignaturas creadas,
    omitidas por error, y detalle de los errores de validación por fila.

    % RF-4.3 ----------------------------------------------------------------
    \item \textbf{Modificación de asignatura.} \label{rf:subjects:update}
    El sistema debe permitir a un gestor o al coordinador de la asignatura modificar sus
    datos.

    \textbf{Entrada:} petición PATCH al recurso de la asignatura con los campos a modificar
    (nombre, código, cuatrimestre, coordinador).

    \textbf{Proceso:} el sistema valida los datos actualizados. Si se cambia el coordinador,
    se actualiza el campo \texttt{coordinator} en la asignatura. El antiguo coordinador
    mantiene su \texttt{SubjectMembership} pero su rol se puede ajustar. Al nuevo
    coordinador se le otorga el rol \texttt{COORDINATOR} en \texttt{SubjectMembership},
    creándola si no existía.

    \textbf{Salida:} respuesta con los datos actualizados de la asignatura.

    % RF-4.4 ----------------------------------------------------------------
    \item \textbf{Eliminación de asignatura.} \label{rf:subjects:delete}
    El sistema debe permitir a un gestor eliminar una asignatura que no tenga exámenes
    asociados.

    \textbf{Entrada:} petición DELETE al recurso de la asignatura.

    \textbf{Proceso:} el sistema verifica que la asignatura no tiene exámenes creados. Si
    los tiene, la eliminación se rechaza para preservar la integridad de los datos. Si no
    tiene exámenes, se eliminan las \texttt{SubjectMembership} asociadas y se elimina la
    asignatura.

    \textbf{Salida:} respuesta con código \acs{http}~204 si se elimina correctamente, o
    código \acs{http}~409 si tiene exámenes asociados.

    % RF-4.5 ----------------------------------------------------------------
    \item \textbf{Gestión de grupos.} \label{rf:subjects:groups}
    El sistema debe permitir al coordinador o al gestor crear, modificar y eliminar grupos
    dentro de una asignatura.

    \textbf{Entrada:} peticiones POST, PATCH o DELETE al recurso de grupos de la
    asignatura, con el identificador del grupo (cadena alfanumérica, por ejemplo
    <<2121>>).

    \textbf{Proceso:} al crear un grupo, se valida la unicidad del identificador dentro de
    la asignatura. Al eliminar un grupo, los estudiantes asignados a él quedan en la
    asignatura sin grupo asignado (mantienen su \texttt{SubjectMembership} con el campo de
    grupo vacío). Un estudiante sin grupo no puede ser convocado a exámenes hasta que se le
    reasigne uno.

    \textbf{Salida:} respuesta con los datos del grupo creado o modificado, o código
    \acs{http}~204 en caso de eliminación.

    % RF-4.6 ----------------------------------------------------------------
    \item \textbf{Asignación de usuarios a asignatura.} \label{rf:subjects:assign}
    El sistema debe permitir al coordinador o al gestor vincular usuarios a una asignatura
    con un rol específico (profesor o estudiante) y asignarlos a un grupo.

    \textbf{Entrada:} petición POST al recurso de miembros de la asignatura con el
    identificador del usuario, el rol (\texttt{TEACHER} o \texttt{STUDENT}) y,
    opcionalmente, el identificador del grupo.

    \textbf{Proceso:} el sistema valida que el usuario exista, esté activo y pertenezca a la
    organización. Se crea una \texttt{SubjectMembership} con el rol y los permisos
    contextuales predeterminados según la configuración de la organización (véase
    \cref{rf:config:default_perms}). Si el rol es \texttt{STUDENT}, se verifica que el
    grupo indicado exista dentro de la asignatura y que el estudiante no esté ya asignado a
    otro grupo en la misma asignatura. Un estudiante puede estar temporalmente sin grupo
    asignado. Si el rol es \texttt{TEACHER}, la asignación a grupo es opcional.

    \textbf{Salida:} respuesta con los datos de la asignación (usuario, rol, grupo,
    permisos).

    % RF-4.7 ----------------------------------------------------------------
    \item \textbf{Desasignación de usuarios de asignatura.} \label{rf:subjects:unassign}
    El sistema debe permitir al coordinador o al gestor desvincular a un usuario de una
    asignatura.

    \textbf{Entrada:} petición DELETE al recurso del miembro de la asignatura
    (\texttt{SubjectMembership}).

    \textbf{Proceso:} si el usuario tiene instancias de examen asociadas como corrector o
    estudiante, el sistema marca la \texttt{SubjectMembership} como
    \texttt{is\_active=False} (desactivación lógica); en caso contrario, se elimina
    físicamente.

    \textbf{Salida:} respuesta con código \acs{http}~204.

    % RF-4.8 ----------------------------------------------------------------
    \item \textbf{Listado de asignaturas del usuario.} \label{rf:subjects:own_list}
    El sistema debe permitir a cualquier usuario autenticado consultar las asignaturas del
    curso activo en las que participa.

    \textbf{Entrada:} petición GET al recurso de asignaturas propias.

    \textbf{Proceso:} el sistema recupera las asignaturas del curso activo de la
    organización del usuario en las que tiene una \texttt{SubjectMembership} activa
    (\texttt{is\_active=True}), incluyendo el rol y el grupo si procede.

    \textbf{Salida:} respuesta con la lista de asignaturas, cada una con su nombre, código,
    cuatrimestre, rol del usuario y grupo asignado.

    % RF-4.9 ----------------------------------------------------------------
    \item \textbf{Consulta detallada de asignatura.} \label{rf:subjects:detail}
    El sistema debe permitir consultar la información completa de una asignatura, incluyendo
    miembros, grupos y exámenes.

    \textbf{Entrada:} petición GET al recurso de la asignatura.

    \textbf{Proceso:} el sistema recupera los datos de la asignatura según el rol del
    usuario que consulta. Los gestores, el coordinador y los profesores obtienen la
    información completa (todos los miembros, grupos y exámenes). Los estudiantes obtienen
    la información básica de la asignatura, su grupo y los exámenes publicados, sin lista
    de estudiantes.

    \textbf{Salida:} respuesta con los datos de la asignatura adaptados al nivel de acceso
    del usuario.

    % RF-4.10 ---------------------------------------------------------------
    \item \textbf{Exportación de asignaturas.} \label{rf:subjects:export}
    El sistema debe permitir a un gestor exportar las asignaturas del curso activo de su
    organización en formato estructurado.

    \textbf{Entrada:} petición GET al recurso de exportación de asignaturas con parámetro
    de formato (\acs{csv} o \acs{json}).

    \textbf{Proceso:} el sistema genera un fichero con los datos de las asignaturas,
    incluyendo las asignaciones de usuarios y grupos, en el formato solicitado.

    \textbf{Salida:} respuesta con el fichero generado como descarga.

\end{functional}
